% =============================================
% CAPÍTULO 9 - PROPOSTA DE DESENVOLVIMENTO
% =============================================
\chapter{PROPOSTA DE DESENVOLVIMENTO}

Para que o projeto seja desenvolvido de forma eficiente e atinja os padrões de qualidade e confiabilidade esperados, contamos com a utilização de diversas tecnologias e práticas de engenharia de software que apoiam o processo de desenvolvimento. A arquitetura adotada segue o padrão de microsserviços orientados a eventos, com Apache Kafka como barramento central de mensagens, permitindo o desacoplamento dos serviços, resiliência a picos de carga e escalabilidade horizontal.

Segundo Newman (2022), microsserviços são uma abordagem para o desenvolvimento de sistemas onde uma única aplicação é construída como um conjunto de serviços pequenos e independentes, cada um executando em seu próprio processo e se comunicando por meio de mecanismos leves. Essa abordagem favorece a implantação independente, a resiliência e a evolução contínua do produto.

O processo de desenvolvimento seguirá metodologias ágeis e práticas como o TDD (\textit{Test-Driven Development}) para garantir a confiabilidade e a qualidade das entregas ao longo do ciclo de vida do projeto.

\section{Tecnologias que serão utilizadas}

O Quadro~\ref{qua:tecnologias} apresenta as tecnologias que serão utilizadas para que o projeto seja concluído, desde sua elaboração até a implementação e operação em produção.

{\small
\begin{longtable}{|p{3.2cm}|p{2.8cm}|p{6.7cm}|p{1.5cm}|}
\caption{Tecnologias Utilizadas} \label{qua:tecnologias} \\
\hline
\textbf{Nome} & \textbf{Fornecedor} & \textbf{Descrição} & \textbf{Versão} \\ \hline
\endfirsthead
\hline
\textbf{Nome} & \textbf{Fornecedor} & \textbf{Descrição} & \textbf{Versão} \\ \hline
\endhead
\hline
\multicolumn{4}{r}{\fontsize{10pt}{10pt}\selectfont\textit{Continua na próxima página}} \\
\endfoot
\hline
\multicolumn{4}{c}{\fontsize{10pt}{10pt}\selectfont \textbf{Fonte:} Elaborado pelos autores (2026).} \\
\endlastfoot
Git & Open Source & Sistema de controle de versão usado para rastrear alterações no código-fonte. & 2.x \\ \hline
GitHub & Microsoft & Plataforma de hospedagem de código-fonte, colaboração e CI/CD via GitHub Actions. & -- \\ \hline
Node.js & OpenJS Foundation & Ambiente de execução JavaScript assíncrono e orientado a eventos, utilizado como runtime de todos os microsserviços. & 20 LTS \\ \hline
TypeScript & Microsoft & Superset de JavaScript com tipagem estática, utilizado em todos os serviços para maior segurança e manutenibilidade. & 5.x \\ \hline
Fastify & Open Source & Framework HTTP para Node.js de alta performance, utilizado como base dos microsserviços. & 5.x \\ \hline
Apache Kafka & Apache / Upstash & Barramento de eventos distribuído, responsável pela comunicação assíncrona entre os microsserviços. & -- \\ \hline
Redis & Upstash & Armazenamento de dados em memória, utilizado para cache de sessões, tokens temporários e rate limiting. & 7.x \\ \hline
PostgreSQL & Open Source & Banco de dados relacional, utilizado para dados estruturados, audit logs e consentimentos LGPD. & 16 \\ \hline
pgcrypto & Open Source & Extensão do PostgreSQL para criptografia de colunas sensíveis com AES-256. & built-in \\ \hline
pgvector & Open Source & Extensão do PostgreSQL para armazenamento e busca de vetores de embeddings, utilizada no pipeline RAG. & 0.7+ \\ \hline
MongoDB & MongoDB, Inc. & Banco de dados documental, utilizado para armazenar transações financeiras e histórico de conversas com criptografia em nível de campo. & 7.x \\ \hline
BullMQ & Open Source & Biblioteca de filas de jobs baseada em Redis, responsável pelo agendamento de sincronizações e notificações periódicas. & 5.x \\ \hline
KafkaJS & Open Source & Cliente Kafka nativo para Node.js, utilizado pelos consumers e producers dos microsserviços. & 2.x \\ \hline
Chart.js + @napi-rs/canvas & Open Source & Biblioteca de renderização de gráficos server-side, responsável pela geração de imagens PNG para envio via WhatsApp. & 4.x \\ \hline
Docker + Docker Compose & Docker, Inc. & Plataforma de conteinerização e orquestração dos serviços em ambiente de desenvolvimento e produção. & 27+ \\ \hline
Nginx & Open Source & Reverse proxy responsável pela terminação TLS e roteamento de requisições externas para o API Gateway. & -- \\ \hline
Let's Encrypt & ISRG & Autoridade certificadora gratuita para emissão e renovação automática de certificados SSL/TLS. & -- \\ \hline
Helmet.js & Open Source & Conjunto de middlewares de segurança HTTP para Node.js (HSTS, X-Frame-Options, X-Content-Type-Options). & 8.x \\ \hline
Zod & Open Source & Biblioteca de validação e parsing de schemas TypeScript, aplicada na camada de produção de eventos Kafka. & 3.x \\ \hline
Drizzle ORM & Open Source & ORM TypeScript para PostgreSQL, utilizado nas operações relacionais com tipagem segura. & -- \\ \hline
Mongoose & Open Source & ODM para MongoDB em Node.js, utilizado no acesso aos dados de transações e histórico. & 8.x \\ \hline
OpenAI API & OpenAI & API de modelos de linguagem (GPT-4o-mini) e embeddings (text-embedding-3-small), utilizada no pipeline conversacional e RAG. & -- \\ \hline
Meta Cloud API & Meta & API oficial da plataforma WhatsApp Business para recepção e envio de mensagens. & v21+ \\ \hline
Pluggy & Pluggy & API de integração com o ecossistema Open Finance do Banco Central do Brasil. & v1 \\ \hline
Grafana Cloud & Grafana Labs & Plataforma de monitoramento, métricas (Prometheus) e alertas operacionais. & -- \\ \hline
\end{longtable}
}

\section{Metodologia de padrões de software}

Os padrões de software são fundamentais para garantir a qualidade, a coerência e a manutenibilidade de sistemas complexos. Segundo Bass, Clements e Kazman (2022), decisões arquiteturais determinam as propriedades de qualidade de um sistema e devem ser tomadas de forma explícita e documentada, uma vez que influenciam diretamente aspectos como desempenho, segurança, disponibilidade e modificabilidade. O Quadro~\ref{qua:padroes} apresenta os padrões adotados neste projeto.

{\small
\begin{longtable}{|p{4.5cm}|p{10cm}|}
\caption{Padrões de Projeto} \label{qua:padroes} \\
\hline
\textbf{Nome} & \textbf{Descrição} \\ \hline
\endfirsthead
\hline
\textbf{Nome} & \textbf{Descrição} \\ \hline
\endhead
\hline
\multicolumn{2}{r}{\fontsize{10pt}{10pt}\selectfont\textit{Continua na próxima página}} \\
\endfoot
\hline
\multicolumn{2}{c}{\fontsize{10pt}{10pt}\selectfont \textbf{Fonte:} Elaborado pelos autores (2026).} \\
\endlastfoot
Metodologia Scrum & Framework ágil para gestão de projetos com foco em entregas frequentes e adaptação contínua. Baseado em ciclos de trabalho chamados \textit{sprints}, com cerimônias de planejamento, revisão e retrospectiva. \\ \hline
Arquitetura de Microsserviços Orientada a Eventos & Os serviços do sistema (API Gateway, Chat Service, Finance Service, Notification Service, Outbound Worker, Audit Service) são desacoplados e se comunicam exclusivamente por meio de eventos no Apache Kafka. Cada serviço possui responsabilidade única e pode ser escalado, implantado e testado de forma independente. \\ \hline
Padrão MSC (\textit{Model-Service-Controller}) & Padrão arquitetural que separa responsabilidades em três camadas: \textit{Controller} (recebe e valida a entrada), \textit{Service} (executa a lógica de negócio) e \textit{Model} (representa e persiste os dados). Aplicado dentro de cada microsserviço. \\ \hline
REST & Estilo arquitetural para APIs internas baseado nos métodos HTTP (GET, POST, PUT, DELETE) e representações de recursos em JSON. Utilizado nas interfaces administrativas e na comunicação síncrona pontual entre serviços. \\ \hline
RAG (\textit{Retrieval-Augmented Generation}) & Padrão de IA que combina busca semântica em base vetorial (pgvector) com geração de linguagem natural pelo LLM. O contexto financeiro relevante do usuário é recuperado e inserido no prompt, tornando as respostas personalizadas e fundamentadas nos dados reais do titular. \\ \hline
TDD (\textit{Test-Driven Development}) & Metodologia de desenvolvimento em que os testes automatizados são escritos antes do código funcional. O ciclo envolve escrever um teste que falha, implementar o código mínimo para fazê-lo passar e refatorar. Garante cobertura de testes desde o início do desenvolvimento. \\ \hline
\end{longtable}
}

\section{Metodologia para testes de software}

Na abordagem escolhida para testes de software, adotaremos o \textit{Test-Driven Development} (TDD). Segundo Richardson (2019), em arquiteturas de microsserviços, testar cada serviço de forma isolada e em conjunto com os demais é essencial para detectar falhas de integração precocemente, especialmente em cenários assíncronos onde a comunicação ocorre por meio de eventos. Para isso, cada método implementado será precedido pela criação de testes unitários que validem o comportamento esperado da unidade de código.

Além dos testes unitários, serão realizados testes de integração para verificar o fluxo completo entre os serviços. Esses testes cobrem cenários como a produção e o consumo de eventos Kafka, o armazenamento e a recuperação de dados no PostgreSQL e no MongoDB, a expiração e a renovação do token de privilégio no Redis e a execução dos jobs agendados via BullMQ. Conforme recomendado por Newman (2022), contratos de eventos entre serviços produtores e consumidores devem ser testados explicitamente para evitar que mudanças em um serviço quebrem silenciosamente os demais.

Para garantir a integridade das funcionalidades de segurança e conformidade com a LGPD, serão realizados testes manuais por meio de roteiros que simulam os fluxos de cadastro, autenticação via OTP, vinculação de contas Open Finance, solicitação de exportação de dados e exercício do direito de eliminação. Esses roteiros cobrem tanto o caminho feliz quanto cenários de erro, expiração de tokens e tentativas de acesso não autorizado.

\section{Métricas de qualidade}

Segundo Bass, Clements e Kazman (2022), atributos de qualidade são propriedades mensuráveis de um sistema que indicam quão bem ele satisfaz as necessidades de seus \textit{stakeholders}. Em sistemas distribuídos orientados a eventos, atributos como disponibilidade, desempenho e segurança devem ser projetados explicitamente na arquitetura, não apenas verificados após a implementação. As métricas de qualidade serão utilizadas para comparar e avaliar o software ao longo do seu ciclo de vida. O Quadro~\ref{qua:metricas} apresenta as métricas adotadas neste projeto.

{\small
\begin{longtable}{|p{4cm}|p{10.5cm}|}
\caption{Métricas de Qualidade} \label{qua:metricas} \\
\hline
\textbf{Fator de Qualidade} & \textbf{Objetivo} \\ \hline
\endfirsthead
\hline
\textbf{Fator de Qualidade} & \textbf{Objetivo} \\ \hline
\endhead
\hline
\multicolumn{2}{r}{\fontsize{10pt}{10pt}\selectfont\textit{Continua na próxima página}} \\
\endfoot
\hline
\multicolumn{2}{c}{\fontsize{10pt}{10pt}\selectfont \textbf{Fonte:} Elaborado pelos autores (2026).} \\
\endlastfoot
Manutenibilidade &
\begin{minipage}[t]{\linewidth}
\begin{itemize}[noitemsep,topsep=2pt,leftmargin=*]
    \item Organização do código em microsserviços com responsabilidade única e separação por camadas (MSC).
    \item Uso de TypeScript com tipagem estrita e validação de schemas via Zod para reduzir erros em tempo de desenvolvimento.
    \item Eventos Kafka com schemas versionados para garantir compatibilidade entre produtores e consumidores.
\end{itemize}
\end{minipage} \\ \hline
Desempenho &
\begin{minipage}[t]{\linewidth}
\begin{itemize}[noitemsep,topsep=2pt,leftmargin=*]
    \item Processamento assíncrono via Kafka deve garantir latência de resposta ao webhook da Meta inferior a 200ms (retorno imediato do HTTP 200).
    \item Pipeline RAG completo (busca semântica + chamada LLM) deve ser concluído em até 30 segundos.
    \item Geração e envio de resumos periódicos para 1.000 usuários deve ser concluído dentro da janela de envio configurada.
\end{itemize}
\end{minipage} \\ \hline
Disponibilidade e Resiliência &
\begin{minipage}[t]{\linewidth}
\begin{itemize}[noitemsep,topsep=2pt,leftmargin=*]
    \item Circuit breakers no Outbound Worker devem isolar falhas da Meta Cloud API sem propagar indisponibilidade para os demais serviços.
    \item Retry policies com backoff exponencial em todas as integrações externas (Pluggy, Meta, OpenAI).
    \item Dead letter queues para mensagens que falham repetidamente, permitindo reprocessamento manual.
    \item Healthchecks em todos os containers Docker com restart automático em caso de falha.
\end{itemize}
\end{minipage} \\ \hline
Segurança e Conformidade LGPD &
\begin{minipage}[t]{\linewidth}
\begin{itemize}[noitemsep,topsep=2pt,leftmargin=*]
    \item Dados financeiros e PII criptografados em repouso (AES-256-GCM) e em trânsito (TLS 1.2+).
    \item Token de privilégio com TTL de 1 hora para acesso a valores monetários, com mascaramento automático após expiração.
    \item Audit log imutável com hash SHA-256 encadeado registrando toda operação sobre dados pessoais.
    \item Suporte completo aos direitos do titular previstos no Art. 18 da LGPD (acesso, portabilidade, eliminação, revogação de consentimento).
\end{itemize}
\end{minipage} \\ \hline
Portabilidade &
\begin{minipage}[t]{\linewidth}
\begin{itemize}[noitemsep,topsep=2pt,leftmargin=*]
    \item Todos os serviços são conteinerizados via Docker, garantindo comportamento idêntico em ambiente de desenvolvimento e produção.
    \item Configuração externalizada via variáveis de ambiente, sem dependência de configurações fixas no código.
\end{itemize}
\end{minipage} \\ \hline
Usabilidade &
\begin{minipage}[t]{\linewidth}
\begin{itemize}[noitemsep,topsep=2pt,leftmargin=*]
    \item Interação exclusivamente via WhatsApp, canal já familiar ao usuário, sem necessidade de instalação de aplicativo adicional.
    \item Respostas do assistente devem ser claras, objetivas e acompanhadas de representações visuais (gráficos em PNG) quando relevante.
    \item Comandos de gerenciamento de dados e consentimentos devem ser executáveis em linguagem natural dentro do próprio chat.
\end{itemize}
\end{minipage} \\ \hline
Observabilidade &
\begin{minipage}[t]{\linewidth}
\begin{itemize}[noitemsep,topsep=2pt,leftmargin=*]
    \item Métricas de cada microsserviço coletadas via Prometheus e visualizadas no Grafana Cloud.
    \item Logs estruturados em JSON com correlação por \texttt{event\_id} para rastreabilidade de fluxos assíncronos entre serviços.
    \item Alertas automáticos para taxa de erro acima do limiar configurado, latência elevada e falhas em jobs agendados.
\end{itemize}
\end{minipage} \\ \hline
\end{longtable}
}
